
% 04. Equilibrium Trapped, Linear Motion Failed.tex

\section{갇힌 균형, 실패한 직선운동}
우리는 아주 어릴 때부터 우주에 두 가지 종류의 근본적인 운동이 존재한다고 배워왔다. 자동차가 도로를 달리고 총알이 날아가는 '직선운동'과, 행성이 태양을 돌고 팽이가 제자리에서 도는 '원운동(회전운동)'이다. 고전 물리학 역시 이 둘을 철저히 분리하여, 직선운동은 선속도와 운동량 보존으로, 회전운동은 각속도와 각운동량 보존으로 설명하며 서로 다른 수식과 잣대를 들이댄다. 우리는 너무도 자연스럽게 이 둘을 전혀 다른 현상으로, 전혀 다른 궤적으로 분리하여 받아들여 왔다.
그러나 디스턴스(DISTANCE)의 서늘한 렌즈를 들이대면, 이 견고한 이분법은 단숨에 무너져 내린다. 과연 이 둘은 정말로 다른 운동일까?

우리는 항상 운동을 '독립된 물체'의 관점에서 바라보는 오류를 범한다. 텅 빈 3차원의 시공간 무대 위에서 저 물체가 이동하고 저 행성이 돈다고 믿는 것이다. 하지만 우주의 실제 하드웨어인 '에너지 갭'의 관점에서 우주를 바라보면, 원자도, 세포도, 인간도, 그리고 은하조차도 텅 빈 공간을 가로지르는 독립된 실체가 아니다. 그들은 모두 맹렬한 모멘텀의 흐름 속에서 내부적으로 팽팽한 구속력을 유지하고 있는 임시적인 '에너지 아일랜드(섬)'들이다.
이 관점에서 보면 운동이란 시공간 상의 위치 변화가 결코 아니다. 우주 전체를 아우르는 절대적 기준점이 존재하지 않는 이상, 모든 운동의 진짜 실체는 오직 에너지 아일랜드들 사이의 '역학적 디스턴스(관계의 거리) 변화'일 뿐이다. 즉 "무엇이 어디로 움직이는가"가 아니라, "어떤 에너지 관계(갭)가 해소되고 있는가"를 묻는 것이 유일한 우주적 진실이다.

이제 어떤 국소적 에너지 아일랜드에, 다른 아일랜드와의 에너지 갭을 해소하려는 거대한 방향성(모멘텀)이 작용했다고 가정해 보자. 만약 그 모멘텀이 아일랜드 전체에 완벽하게 균일하게 전달되어 저항 없이 뻗어 나간다면, 우리는 그것을 '직선운동'이라 부른다. 하지만 실제 우주는 그렇게 단순하지 않다. 팽이나 행성, 은하는 맹렬한 모멘텀을 받더라도 쉽게 쪼개지거나 분열되지 않는다. 그들 내부에 스스로를 하나의 아일랜드로 묶어두려는 팽팽한 '내부 구속력'이 존재하기 때문이다.
자신의 디스턴스를 해소하기 위해 맹렬하게 바깥으로 뻗어 나가며 분리되려는 모멘텀(기존 물리학의 원심력)과, 아일랜드를 하나의 조직으로 꽉 쥐고 놓지 않으려는 내부 구속력(구심력)이 한 치의 양보 없이 정면으로 충돌할 때 어떤 일이 벌어지는가? 뻗어 나가려던 직선의 모멘텀은 구속력을 이기지 못하고 꺾여 들어가며, 결국 조직 내부의 순환하는 궤도로 갇혀버리게 된다. 디스턴스 에세이는 장엄하게 선언한다. 회전(원운동)은 결코 직선과 독립된 별개의 운동이 아니다. 회전은 구속력에 의해 뻗어 나가지 못하고 '실패해 버린 직선운동', 즉 영원히 갇혀버린 맹렬한 '갇힌 균형(Trapped Equilibrium)'일 뿐이다.

결국 직선과 원은 서로 반대되는 궤적이 아니다. 디스턴스(에너지 갭)를 해소하려는 완벽하게 동일한 단 하나의 모멘텀이, 구속력의 저항 없이 자유롭게 뻗어 나가면 직선이 되고, 내부 연결망에 붙잡히면 원이 되어 돌 뿐이다. 직선과 원은 같은 길을 전혀 다른 형태로 걷고 있는 완벽한 동형(Isomorphism)인 것이다.
이 파괴적인 통찰은 우리가 '관성(Inertia)'이라 부르는 우상의 의미마저 뿌리째 전복해 버린다. 물리학은 관성을 '물체가 현재 상태를 고집스레 유지하려는 의지나 성질'로 정의해 왔다. 그러나 디스턴스의 렌즈로 보면 관성은 성질이 아니다. 관성의 진정한 물리적 실체는 '단절(Severance)'이다. 특정 스케일에서 주변의 어떤 우주적 존재와도 새로운 에너지 소통(모멘텀)을 형성하지 못하고 뚝 끊어져 버린 '관계의 부재' 상태인 것이다. 따라서 우주 공간에 가만히 멈춰 있는 돌덩이와, 시속 10만 킬로미터로 등속 비행하는 우주선은 시공간의 잣대(속도)로는 극단적으로 달라 보이지만, 에너지 교류가 완벽히 끊어져 있다는 디스턴스의 층위에서는 소름 돋도록 완벽하게 동일한 상태이다.

그렇다면, 이 우주를 묶고 있는 가장 거대한 힘인 '중력(Gravity)'은 무엇인가? 직선과 원이 본질적으로 동일하다면, 사과를 땅으로 떨어뜨리고 달을 지구에 붙잡아 두는 이 현상의 진짜 실체는 무엇인가? 중력은 텅 빈 공간을 가로지르는 마법의 밧줄도, 아인슈타인이 말한 3차원 시공간의 기하학적 휘어짐(곡률)도 결코 아니다. 우주에 스스로 휘어지는 무대 따위는 존재하지 않는다. 중력은 거시적 아일랜드와 국소적 아일랜드 사이에 존재하는 거대한 에너지 갭을 소거하고, 끝내 융합하여 '단 하나의 거대한 아일랜드'가 되기 위해 맹렬하게 뻗어 나가는 우주적 모멘텀의 가장 노골적인 역학적 표출일 뿐이다.

결국 이 거대한 우주에는 직선운동과 원운동이라는 두 종류의 운동이 존재하지 않는다. 중력이라는 분리된 힘도, 공간의 팽창도, 시공간의 흐름도 진짜 하드웨어가 아니다. 우주에는 오직 단 하나, 가장 거대한 에너지 갭이 가장 완벽한 평활화(정지)를 향해 스스로의 모멘텀을 처절하게 무너뜨리며 섞어 내는 '단일한 추락의 과정'만이 영구히 지속되고 있을 뿐이다.
